%%%%%%%%%%%%%%%%%%%%%%%%%%%%%%%%%%%%%%%%%%%%%%%%%%%%%%%%%%%%%%%%%%%%%
%%%%%%%%%%%%%%%%%%%%%%%%%%%%%%%%%%%%%%%%%%%%%%%%%%%%%%%%%%%%%%%%%%%%%
%%%%%%%%%%%%%%%%%%%%%%%%%%%%%%%%%%%%%%%%%%%%%%%%%%%%%%%%%%%%%%%%%%%%%
\section{Introduction}\label{sec:introduction}

In their work on real Morsifications and quiver mutations, S.~Fomin, P.~Pylyavskyy, E.~Shustin and D.~Thurston put forth the following conjecture:\\
\vspace{-0.2cm}
\begin{fthm2}[\cite{FPST}]\label{conj:FPST}
Given two real morsifications of real isolated plane curve singularities, the following are equivalent:
\begin{itemize}
\item[$(i)$] the two singularities have the same complex topological type;
\item[$(ii)$] the quivers associated with the two morsifications are mutation equivalent.
\end{itemize}
\end{fthm2}
\vspace{0.2cm}
\noindent The precise reference is \cite[Conjecture 5.5]{FPST}. In my view, the Main Conjecture is remarkable, as it asserts that isolated plane curve singularities are
topologically classified by the mutation classes of their associated quivers. The former are central objects of study in algebraic geometry and the theory of singularities. The study of plane curves and their singularities has had a prominent role in mathematics for centuries, see e.g.~\cite{brieskorn1986, casas2000, ghys2017, milnor1968, wall2004, zariski2006}. The latter, the study of quivers and their mutation, is a much more recent enterprise, particularly bolstered by the series of articles \cite{Berenstein2005,Fomin2001,Fomin2003} and the subsequent study of cluster algebras. In a sense, $(i)$ and $(ii)$ in the Main Conjecture belong to apparently different areas of mathematics: topology and algebraic geometry, for $(i)$, and combinatorics and graph theory, for $(ii)$.\\

Heretofore, both directions of the equivalence in the Main Conjecture are open. The case of simple singularities, equivalently quivers of finite type, is established in \cite[Theorem 18.3]{FPST}. See also \cite[Conjecture 7.11 \& Corollary 9.10]{FPST}. The object of this article is to establish the equivalence $(ii)\Longrightarrow(i)$ of the Main Conjecture in a central case. Namely, when the quivers are associated to irreducible malleable divides. It is a central case because of the likelihood of \cite[Conjecture 14.10]{FPST}, which asserts that all algebraic divides are malleable. See \Cref{fig:divides-lissajous} for an instance of a divide $D$ and its quiver $Q(D)$. More generally, our main result establishes that the quiver mutation class of any malleable divide, irreducible or not, uniquely recovers the integral monodromy module of the singularity.\\


The precise main result of the article reads as follows:
\vspace{0.1cm}
\begin{fthm}\label{thm:main} Let \((C,0)\) be an isolated plane curve singularity, $D$ a divide arising from a real Morsification of $(C,0)$ and $Q(D)$ its quiver. Suppose that $D$ is malleable. Then the quiver mutation class $[Q(D)]$ uniquely determines the integral monodromy module of $(C,0)$.
\end{fthm}
\vspace{0.2cm}
In particular, the Main Theorem proves the implication $(ii)\Longrightarrow(i)$ of the ``Main Conjecture'' \cite[Conjecture 5.5]{FPST} for irreducible plane curve singularities with malleable divides. Indeed, the integral monodromy module of $(C,0)$ determines the Alexander polynomial $\Delta_C(t)$, which itself determines the topological type of an irreducible plane curve singularity. There are also infinitely many instances of reducible plane curve singularities which are determined by the integral monodromy module, or even its one-variable Alexander polynomial, cf.~\Cref{sec:illustrative_examples}, and thus the Main Theorem can be applied more generally to resolve $(ii)\Longrightarrow(i)$ of the ``Main Conjecture'' in such cases.\\

Note that the Main Theorem can also be phrased as follows. Let \((C_1,0)\) and \((C_2,0)\) be two isolated plane curve singularities, with $D_1,D_2$ two divides arising from two real Morsifications of these two singularities. Suppose that $D_1$ and $D_2$ are malleable and that the integral monodromy module of either \((C_1,0)\) or \((C_2,0)\) determines its complex topological type. Then\\
\begin{center}
\fbox{$Q(D_1)\sim_{\mathrm{mut}}Q(D_2)\Longrightarrow (C_1,0) \mbox{ and }(C_2,0) \mbox{ have the same complex topological type.}$}
\end{center}
\vspace{0.2cm}

\begin{figure}[htbp]
\centering

\scalebox{2}{\LissajousFiveFourQuiverPicture}

\caption{An irreducible malleable divide $D$ for the singularity $(C,0)=\{x^4=y^5\}$, in black, together with its
associated quiver $Q(D)$, with blue vertices and red arrows. The quiver $Q(D)$ is the quiver associated to the Grassmannian $\mbox{Gr}(4,9)$, cf.~\cite[Remark 16.2]{FPST}. The link of this irreducible singularity is the $(4,5)$-torus knot.}
\label{fig:divides-lissajous}
\end{figure}

A core challenge of the implication $(ii)\Longrightarrow (i)$ of the Main Conjecture, which the Main Theorem overcomes, is as follows. Given a divide $D$ of a real Morsification of an isolated plane curve singularity $(C,0)$, its quiver $Q(D)$ has a direct connection to the geometry of $(C,0)$. To wit, the vertices of $Q(D)$ can be identified with the vanishing cycles associated to the Morsification, and the arrows of $Q(D)$ exactly capture the intersection numbers between such vanishing cycles. In fact, \cite[Theorem 4.4]{FPST} shows that $Q(D)$ determines the topological type of the singularity. The challenge is that another quiver $\mu(Q)\in[Q(D)]$ in the quiver mutation class of $Q(D)$ will, a priori, have little to do with the geometry of the singularity $(C,0)$. Among other issues: such quiver $\mu(Q)$ will almost never be associated to a divide, nor there will be a compatible $\{\oplus,\ominus,\bullet\}$-coloring, nor the arrows of $\mu(Q)$ will have any direct meaning in relation to $(C,0)$. Plainly, suppose we asked a computer to mutate the quiver $Q(D)$, say 8937284651872 times, and it would return the resulting quiver $\mu(Q)$ to us, but without specifying the mutation sequence it used. How would we recover from such output $\mu(Q)$ any information such as the characteristic polynomial of the monodromy of the singularity $(C,0)$?\\

In order to prove the Main Theorem, we develop the study of a 3-Calabi-Yau differential bigraded Ginzburg algebra $\Gamma(Q,W,d)$ associated to a graded quiver with potential $(Q,W,d)$. Such a dg algebra incorporates an Adams grading associated to an arrow-grading $d:Q_1\lr\Z$ for which the potential $W$ is homogeneous of degree one. The two key new ideas to prove the Main Theorem are:
\begin{itemize}
\item[$(i)$] First, for quivers associated to plabic fences, we show that the cokernel of a certain equivariant Euler pairing in the derived category of continuous finite-dimensional dg modules over $\Gamma(Q,W,d)$, weighted with respect to the internal $d$-grading, is an invariant of the quiver mutation class $[Q]$. Here we choose a non-degenerate potential $W$ and a corresponding arrow-grading $d$, which we show are uniquely associated to such quivers. The content of Sections \ref{sec:graded-QP}, \ref{sec:euler} and \ref{sec:dilation_plabic_fence} is aimed at developing and establishing these results.\\

\item[$(ii)$] Second, we show that such cokernel is isomorphic, as a $\Z[t,t^{-1}]$-module, to the torsion part of the Alexander module of the smooth link associated to the plabic fence. This is achieved in \Cref{sec:polynomial_plabic_fence}, by establishing a combinatorial expression for the equivariant Euler pairing and relating it to a Seifert matrix. Such results, along with $(i)$, are then specialized to algebraic links in \Cref{sec:proof_main} to finally establish the Main Result.
\end{itemize}

The results needed to implement $(i)$ are representation-theoretic in nature, building on the study of quivers with potential, as initiated in \cite{DWZ} by H.~Derksen, J.~Weyman, and A.~Zelevinsky, and the derived categorical enhancements of B.~Keller and D.~Yang in \cite{KellerYang}, and C.~Amiot and S.~Oppermann in \cite{AO}. Part of the new results establishing $(i)$ are developed in Sections \ref{sec:graded-QP}, \ref{sec:euler} and \ref{sec:dilation_plabic_fence} and might be of independent interest. It should also be noted that differential bigraded Ginzburg algebras and equivariant Euler pairings have been explored by L.~Fan, B.~Keller and Y.~Qiu in \cite{FanKellerQiu2024}, by R.~Fujita and K.~Murakami in \cite{FujitaMurakami2022}, and see also \cite{Keller2020Cartan} and recent work of Y.~Wu and L.~Fan on Higgs categories for differential bigraded Ginzburg algebras \cite{Wu2026Graded}.\\

From a combinatorial perspective, we obtain a perhaps surprising invariance result from our representation-theoretic results. Namely, that for any quiver $Q\in[Q(\bG)]$ in the mutation class of a quiver $Q(\bG)$ of a plabic fence $\bG$, there exists a canonical arrow grading $d:Q_1\lr\Z$, unique up to certain equivalence, such that the cokernel of the $|Q_0|\times|Q_0|$ matrix $\Eul(t)$ with entries
\begin{equation}\label{eq:euler-formula-intro}
 \boxed{
 \Eul(t)_{ij}
 :=(1-t)\delta_{ij}
 -\sum_{a:j\to i}t^{d(a)}
 +\sum_{a:i\to j}t^{1-d(a)}}
\end{equation}
is the torsion part of the Alexander module of the smooth link associated to $\bG$. In particular, its determinant is the one-variable Alexander polynomial of such a link. The sums in \eqref{eq:euler-formula-intro} run over the arrows $a\in Q_1$ of the quiver $Q$, and $a:i\to j$ denotes an arrow with source vertex $i\in Q_0$ and target vertex $j\in Q_0$. See~\Cref{ssec:graded_Euler_pairing} and \Cref{thm:unique-fence-G1}. Thus, in the case of a malleable divide $D$ of a singularity $(C,0)$, even if many aspects of a quiver $Q$ in $[Q(D)]$ might not have a direct interpretation in terms of $(C,0)$, the Main Theorem shows that it still remembers the action of the algebraic monodromy of $(C,0)$. A combinatorial proof of such fact remains unknown to the author, as I obtained the expression \eqref{eq:euler-formula-intro} by using the weight decomposition of the Ext groups of the simple modules of a differential bigraded Ginzburg algebra, and its invariance is deduced from rather categorical considerations.\\

\noindent The article concludes with \Cref{sec:illustrative_examples} and \Cref{sec:comments_and_examples}, where several detailed examples and related applications are provided. In particular:

\begin{enumerate}
    \item We provide a counterexample to \cite[Conjecture 6.17]{FPST}, exhibiting two plabic graphs which are not move-and-switch equivalent but have mutation-equivalent quivers. These are distinguished by using the top $a$-degree part of their HOMFLY polynomials.\\

    \item We show that the implication $(ii)\Longrightarrow(i)$ in the Main Conjecture holds for all plane curve singularities with Milnor number $\mu\leq16$. In fact, a table of all such singularities with their associated Alexander polynomials is provided, to illustrate how the Main Result and known invariants suffice to classify all of these 74 singularities via their quiver mutation classes.\\
\end{enumerate}


%%%%%%%%%%%%%%%%%%%%%%%%%%%%%%%%%%%%%%%%%%%%%%%%%%%%%%%%%%%%%%%%%%%%%
%%%%%%%%%%%%%%%%%%%%%%%%%%%%%%%%%%%%%%%%%%%%%%%%%%%%%%%%%%%%%%%%%%%%%
%%%%%%%%%%%%%%%%%%%%%%%%%%%%%%%%%%%%%%%%%%%%%%%%%%%%%%%%%%%%%%%%%%%%%
\noindent{\bf Acknowledgements}. I am very thankful to Sergey Fomin and Bernhard Keller for valuable comments that helped improve the original manuscript. I am grateful to the authors of \cite{FPST} for producing such an enticing conjecture. Finally, I thank PCMI-IAS for its hospitality during the research program ``Knotted Surfaces in Four-Manifolds'', where parts of this article were developed. R.C.~is supported by the National Science Foundation under the grant DMS-2505760, and a UC Davis College of L\&S Dean's Fellowship.\hfill$\Box$

\vspace{0.3cm}
%%%%%%%%%%%%%%%%%%%%%%%%%%%%%%%%%%%%%%%%%%%%%%%%%%%%%%%%%%%%%%%%%%%%%
%%%%%%%%%%%%%%%%%%%%%%%%%%%%%%%%%%%%%%%%%%%%%%%%%%%%%%%%%%%%%%%%%%%%%
%%%%%%%%%%%%%%%%%%%%%%%%%%%%%%%%%%%%%%%%%%%%%%%%%%%%%%%%%%%%%%%%%%%%%
\noindent{\bf Conventions}.
We work over \(\kk=\C\).  Quivers are finite and have no loops unless stated
otherwise.  Path multiplication is right-to-left: if \(a:i\to j\) and
\(b:j\to k\), then \(ba\) means first \(a\), then \(b\). Throughout the manuscript $\Lambda:=\Z[t,t^{-1}]$ will denote the Laurent polynomial ring in the $t$-variable, whose units are $\pm t^{m}$ for any $m\in\Z$. \hfill$\Box$